\section{Implementation}
\label{sec:impl}

\paragraph{Engine} We use SGLang v0.5.10 \cite{zheng2024sglang} \emph{unmodified}, with
paged KV \cite{kwon2023pagedattention} and its hierarchical cache enabled at a host ratio
of $2\times$ with write-through policy. The model is Qwen2.5-Coder-7B-Instruct with YaRN
extrapolation to a 96K context, whose weights and KV cache fit a single 24\,GB RTX 3090 at
a static memory fraction of 0.85, leaving a device pool of 101,502 tokens and a host tier
of 203,004. The engine exposes token usage, running-request counts, and cache statistics on
its metrics endpoint; it has no concept of a session and no notion of suspension. Nothing
in the engine was changed to support this paper---the design lives entirely outside it.

\paragraph{Harness} A session runner drives the replay: for each session it emits an
arrival event, then for each round it waits for a permit, submits the round to the engine,
records the step (prompt, cached and uncached tokens, first-token latency, duration), and
simulates the inter-round tool wait before emitting the next round's arrival. The runner
implements the round-boundary gate: a session that is not in the admit set blocks before
its next request is issued, keeping its process state and its engine-side KV. The runner is
written in Rust and is the only component that talks to the engine.

\paragraph{Controller} The controller is a Python process of 540 lines that polls
the engine's metrics and the runner's event logs every two seconds, maintains the ledger of
\S\ref{sec:m1}, evaluates the valve and the shedding rules of \S\ref{sec:m4}, and rewrites
the permit file atomically with full-state semantics: each cycle publishes the complete
admit set and the complete paused set, so the runner never acts on a stale delta. Decision
records---resident bytes, projection, budget, valve state, running count, shed set---are
appended to a log that the evaluation reads. The controller never opens the trace file: the
only quantities it sees are the ones the deployment produces at run time.

\paragraph{Arms} All three arms use the same engine, model, trace, memory fraction, and
cache configuration; they differ only in the control process, which is selected by a policy
string. The \emph{request-level} arm runs with no controller at all---the engine admits
every round on arrival and its own scheduler arbitrates. The \emph{agent-level} arm runs a
Concur-style AIMD controller \cite{concur} that adjusts a concurrency window from observed
utilisation and a cache-thrash signal, pausing and resuming sessions at step boundaries
with the same named-permit mechanism. The \emph{capacity-accounted} arm is the design of
\S\ref{sec:design}. Using one permit mechanism for both controllers matters: the comparison
then isolates the decision rule (AIMD on a window versus accounting on a ledger), not the
engineering of pausing.

\paragraph{Reproducibility} A run is one command per arm and writes a directory containing
the per-step log, the controller's decisions, the admission events and a metadata file with
the code hashes, trace hash, GPU, and pool size. The evaluation scripts in the artifact
recompute every number in \S\ref{sec:eval} from those directories.
